\chapter{Tableau complet des coûts de la campagne}
Le tableau suivant couvre les 24 problèmes de la campagne, tous avec
$\alpha=1$. La construction est exécutée une fois. Les colonnes locale
et tabou indiquent la moyenne et l'écart-type d'échantillon des trois graines
$0,1,2$. Les valeurs sont exprimées dans les unités de coordonnées de
chaque instance ; elles ne sont pas comparables en valeur brute d'une
instance à l'autre.

Les quatre résolutions PLNE sont présentées séparément au
tableau~\ref{tab:exacts}. Les 20 cas exacts ignorés n'ont aucun coût de
solution à renseigner. Les décimales sont arrondies pour la lecture ;
les JSON conservent la précision de calcul.

\begin{table}[htbp]\centering
{\fontsize{9}{12}\selectfont\setlength{\tabcolsep}{4pt}
\begin{tabular}{lrrrr}
\toprule
Instance & $p$ & Construction & Locale : moyenne $\pm s$ & Tabou : moyenne $\pm s$ \\
\midrule
ulysses16 & 3 & 113{,}31 & 68{,}27 $\pm$ 0{,}00 & 68{,}27 $\pm$ 0{,}00 \\
ulysses16 & 5 & 89{,}19 & 65{,}55 $\pm$ 0{,}15 & 65{,}38 $\pm$ 0{,}00 \\
ulysses16 & 10 & 76{,}44 & 63{,}58 $\pm$ 0{,}39 & 63{,}29 $\pm$ 0{,}02 \\
ulysses22 & 3 & 137{,}72 & 83{,}14 $\pm$ 0{,}00 & 83{,}14 $\pm$ 0{,}00 \\
ulysses22 & 5 & 108{,}79 & 72{,}54 $\pm$ 0{,}83 & 71{,}57 $\pm$ 0{,}00 \\
ulysses22 & 10 & 82{,}36 & 65{,}83 $\pm$ 0{,}25 & 65{,}32 $\pm$ 0{,}07 \\
att48 & 3 & 78\,299{,}77 & 72\,273{,}64 $\pm$ 1\,076{,}65 & 70\,903{,}13 $\pm$ 0{,}00 \\
att48 & 5 & 63\,668{,}04 & 56\,467{,}99 $\pm$ 275{,}98 & 55\,615{,}92 $\pm$ 202{,}51 \\
att48 & 10 & 48\,833{,}89 & 45\,563{,}09 $\pm$ 455{,}55 & 42\,790{,}21 $\pm$ 59{,}20 \\
eil51 & 3 & 859{,}99 & 828{,}26 $\pm$ 0{,}60 & 827{,}57 $\pm$ 0{,}00 \\
eil51 & 5 & 735{,}00 & 708{,}56 $\pm$ 4{,}82 & 698{,}17 $\pm$ 2{,}64 \\
eil51 & 10 & 630{,}47 & 563{,}99 $\pm$ 8{,}93 & 550{,}98 $\pm$ 7{,}44 \\
berlin52 & 3 & 16\,352{,}32 & 14\,095{,}02 $\pm$ 19{,}25 & 14\,090{,}96 $\pm$ 12{,}22 \\
berlin52 & 5 & 14\,916{,}74 & 12\,014{,}92 $\pm$ 148{,}48 & 11\,813{,}49 $\pm$ 22{,}63 \\
berlin52 & 10 & 12\,155{,}57 & 9\,869{,}34 $\pm$ 121{,}56 & 9\,594{,}29 $\pm$ 121{,}71 \\
st70 & 3 & 1\,897{,}37 & 1\,762{,}49 $\pm$ 11{,}55 & 1\,753{,}24 $\pm$ 0{,}00 \\
st70 & 5 & 1\,572{,}53 & 1\,388{,}72 $\pm$ 34{,}01 & 1\,360{,}17 $\pm$ 6{,}86 \\
st70 & 10 & 1\,105{,}82 & 1\,074{,}19 $\pm$ 20{,}63 & 1\,006{,}90 $\pm$ 4{,}32 \\
kroA100 & 3 & 77\,233{,}40 & 76\,618{,}84 $\pm$ 645{,}19 & 75\,202{,}82 $\pm$ 496{,}12 \\
kroA100 & 5 & 68\,427{,}57 & 60\,284{,}12 $\pm$ 1\,653{,}97 & 58\,433{,}18 $\pm$ 59{,}57 \\
kroA100 & 10 & 46\,042{,}25 & 43\,469{,}80 $\pm$ 862{,}27 & 41\,457{,}13 $\pm$ 256{,}17 \\
rd100 & 3 & 30\,724{,}65 & 26\,625{,}74 $\pm$ 394{,}09 & 26\,393{,}33 $\pm$ 17{,}48 \\
rd100 & 5 & 21\,243{,}25 & 20\,528{,}93 $\pm$ 567{,}28 & 19\,541{,}48 $\pm$ 55{,}80 \\
rd100 & 10 & 16\,692{,}42 & 14\,072{,}26 $\pm$ 254{,}50 & 13\,604{,}29 $\pm$ 163{,}40 \\
\bottomrule
\end{tabular}
}
\caption{Coûts moyens de construction, recherche locale et recherche tabou,
sur tous les cas de la campagne de démonstration.}
\sourcetable
\end{table}
\FloatBarrier

La source détaillée originale est \code{resultats/campagne/runs.csv}.
La copie \code{resultats/runs_reconcilies.csv} ajoute uniquement le cas
PLNE ignoré terminal, sans solution ni coût. Le fichier
\code{resultats/provenance_rapport.json} décrit cette opération et conserve
les empreintes des fichiers ayant servi à la rédaction.
